\documentclass[12pt]{article}
\usepackage[a4paper, margin=1in]{geometry}
\usepackage{graphicx}
\usepackage{hyperref}
\usepackage{listings}
\usepackage{xcolor}

\title{WhatsApp Wrapped System\
CS108 Project Report}
\author{Tamizhinian S. and Muhammed Safwan R}
\date{April 29, 2026}

\begin{document}

\maketitle
\tableofcontents

\section{Introduction}

This project implements a WhatsApp Wrapped system that simulates group chat conversations, analyzes user behavior, and presents the results as an interactive slideshow.

The system is divided into three main parts: chat generation, statistical analysis, and frontend visualization. The goal is to model realistic conversation patterns and extract meaningful insights such as activity trends, response behavior, and user engagement.


\section{Project Structure}

The system follows a pipeline-based architecture:

\begin{math}
vocabulary.txt \rightarrow chat.py \rightarrow chat.txt \rightarrow analyze.py \rightarrow data.json \rightarrow app.js
\end{math}

\subsection{Components}

\begin{itemize}
\item \textbf{generator/chat.py}: Generates WhatsApp-style chat data.
\item \textbf{generator/personality.json}: Define personality traits for \texttt{chat.py}
\item \textbf{chat.txt}: Stores generated conversations.
\item \textbf{parser/analyze.py}: Parses chat data and computes statistics.
\item \textbf{data.json}: Computed results in a json file.
\item \textbf{web/app.js}:
\item \textbf{web/index.html}:
\item \textbf{web/style.css}:
\end{itemize}

\section{Chat Generation Design}

\subsection{Personality Configuration}

Instead of hardcoding personalities, we used an external configuration(json) file.

\texttt{personality.json:}

Each user is defined using the following attributes:

\begin{itemize}
\item \textbf{Name}: Identifier of the user
\item \textbf{Active Hours}: Time intervals when the user is likely to send messages
\item \textbf{Yapperness}: Probability of sending messages (higher means more talkative)
\item \textbf{Average Message Length}: Controls the number of words in a message
\item \textbf{Ghoster}: Probability of sending multiple messages consecutively
\end{itemize}

Example:

\begin{verbatim}
{
"Name": "Nishit",
"Active Hours": [[9,12], [18,23]],
"Yapperness": 0.7,
"Average Message Length": 8,
"Ghoster": 0.6
}
\end{verbatim}

\subsection{Message Generation}

idk yap smtg

\section{Parser and Statistics anlaysis}

The parser \texttt{analyze.py} reads the chat line-by-line and extracts structured information using regular expressions using re library. See library section \cite{python}

\subsection{Parser}

Each message is parsed into:
\begin{itemize}
\item Timestamp
\item Sender
\item Tag 
\item Message text
\end{itemize}

Regex used in our code was:
\begin{verbatim}
^\[(\d{2}/\d{2}/\d{4}), (\d{2}:\d{2}:\d{2})\] ([^\[]+) \[(?:REPLY TO\]([^:]+)
|INDEPENDENT\]):\s*(.*)$
\end{verbatim}

\subsection{Statistical Definitions}

Each message in the chat has a timestamp $t$, a sender $s$, and message content. The following definitions describe how each statistic is computed.

\subsubsection{Total Messages}

For each user $u$, we count the number of messages they send:

\subsubsection{Total Words}

For each message, we count the number of words and sum them for each user:


\subsubsection{Night Owl}

We define night-time as messages sent between 12 AM and 4 AM.

$NightCount(u) = \text{number of messages sent by } u \text{ between 0 and 4 hours}$

\subsubsection{Ghost}

A ghost is a user whose messages receive the least replies.

$Replies(u) = \text{number of times messages of } u \text{ are replied to}$

$Ghost = \text{user with minimum } Replies(u)$

\subsubsection{Conversation Starter}

We define the conversation starter using a scoring function instead of a fixed threshold.

Let two consecutive messages have timestamps $t_{i-1}$ and $t_i$. The time gap is:

$Gap_i = {t_i - t_{i-1}}$

We introduce a parameter $T$ (called \textit{time\_range}) which represents the minimum gap unit in hours (in our implementation $T = 2$).

We convert the gap into discrete levels:

$n = \left\lfloor \frac{Gap_i}{T} \right\rfloor$

If $n = 0$, the gap is too small and is ignored. If $n \geq 1$, the message is considered as starting a new conversation after a gap of approximately $nT$ hours.

For each user $u$, we maintain a frequency count:

$freq_u[n] = \text{number of times user } u \text{ sends a message after a gap of level } n$

We then define a score for each user:

$Score(u) = \sum_{n \geq 1} (n \cdot T)^x \cdot (freq_u[n])^y$

where:
\begin{itemize}
\item $x$ controls how much importance is given to larger time gaps
\item $y$ controls how much importance is given to repeated conversation starts
\end{itemize}

In our implementation, we use $x = 1$ and $y = 2$, so the score becomes:

$Score(u) = \sum_{n \geq 1} (n \cdot T) \cdot (freq_u[n])^2$

The conversation starter is defined as:

$ConversationStarter = \text{user with maximum } Score(u)$

\textbf{Interpretation:}
\begin{itemize}
\item Larger gaps (higher $n$) contribute more to the score
\item Users who repeatedly restart conversations are rewarded more (due to squaring)
\item Small gaps are ignored automatically since $n = 0$
\end{itemize}

This formulation directly matches the implementation, where gaps are bucketed into levels and weighted using both gap size and frequency.

\subsubsection{Busiest Day}

We count messages for each date:

$DayCount(d) = \text{number of messages on day } d$

$BusiestDays = \text{top5 day with maximum } DayCount(d)$

\subsubsection{Longest Silence}

We compute gaps between consecutive messages:

$Gap = t_i - t_{i-1}$

$LongestSilence = \max(\text{all such gaps})$

\subsubsection{Most Used Emoji}

We extract emojis from messages and count frequency:

$EmojiCount(e) = \text{number of times emoji } e \text{ appears}$

Top 3 emojis are selected for each user.

\subsubsection{Average Response Time}

For each message, we find the next message sent by a different user:

$ResponseTime = t_{next} - t_{current}$

We ignore gaps greater than 3 hours, as it it might not be a reply

For each user:

$AvgResponse(u) = \text{median of all response times}$

\subsubsection{Hype Person}

The hype person is the fastest responder:

$HypePerson = \text{user with minimum } AvgResponse(u)$

\subsubsection{Activity Concentration (Custom Statistic)}

We define a custom statistic called \textbf{Activity Concentration} to measure how concentrated a user's messaging activity is across different hours of the day.

Let there be $p$ users in the chat. We construct an \textbf{activity matrix}:

$A \in R^{p \times 24}$

where each row corresponds to a user and each column corresponds to an hour of the day (from 0 to 23).

Each entry is defined as:

$A[u, h] = \text{number of messages sent by user } {\bf {u}} \text{ during hour } \bf h$

Thus, for each user $u$, we obtain a 24-dimensional activity vector:

$A_u = [a_0, a_1, a_2, \ldots, a_{23}]$

where $a_h$ is the number of messages sent in hour $h$.

To quantify how concentrated this activity is, we compute the \textbf{Euclidean norm} (L2 norm) of this vector:

$C(u) = \sqrt{a_0^2 + a_1^2 + a_2^2 + \cdots + a*{23}^2}$

In our implementation, this is computed efficiently using NumPy:

\begin{verbatim}
norms = np.linalg.norm(activity, axis=1)
\end{verbatim}

where each row of the matrix \texttt{activity} corresponds to a user's hourly activity vector.

\textbf{Interpretation:}

\begin{itemize}
\item If a user's messages are concentrated in a few hours, some entries of $A_u$ are large, leading to a higher norm $C(u)$.
\item If a user's messages are spread evenly across many hours, the entries are more uniform, resulting in a lower norm.
\end{itemize}

\textbf{Example:}

\begin{itemize}
\item User A: $[0,0,0,10,0,\ldots]$ → high concentration → high norm
\item User B: $[2,2,2,2,2,\ldots]$ → spread activity → lower norm
\end{itemize}

Thus, Activity Concentration captures whether a user is active only during specific time windows or throughout the day, which cannot be inferred from total message count alone.

\textbf {Reason why we Choose this:} I wanted to use the linalg.norm function as it was the time when I did the python last activity. So this how I arrived at this function

\section{Frontend Design}

The frontend is implemented as a full-screen slideshow using HTML, CSS, and JavaScript, with Bootstrap used for layout and responsiveness.

\subsection{Structure}

The interface consists of two main sections:
\begin{itemize}
\item \texttt{#landing}: Displays the title and start button
\item \texttt{#slideshow}: Displays statistics as slides
\end{itemize}

Initially, only the landing screen is visible. On clicking the start button, the slideshow is activated by toggling visibility using CSS classes. 

\subsection{Slideshow Mechanism}

All slides are dynamically rendered inside the \texttt{#stats} container using JavaScript. Each slide replaces the previous content, creating a step-by-step presentation of statistics.

A progress bar at the top visually indicates the current slide position.

\subsection{Data Integration}

The frontend reads data from \texttt{data.json}, which contains:
\begin{itemize}
\item group-level statistics
\item per-user statistics
\end{itemize}

JavaScript processes this data and generates slide content accordingly.

\subsection{Styling and Visualization}

Slides are styled using CSS with gradient backgrounds (e.g., \texttt{.ghost}, \texttt{.hype}, \texttt{.owl}) to differentiate statistics. 

Rank-based tables are used to display top users, with visual emphasis on higher ranks.

\subsection{Animation and Responsiveness}

Slides use a smooth ``slide-up'' animation for transitions. The layout is responsive using Bootstrap and media queries, ensuring proper display on different screen sizes.

\subsection{Chart.js Integration}

In addition to tabular displays, we used \textbf{Chart.js} to visualize certain statistics such as message counts and word counts.

The data from \texttt{data.json} is converted into arrays of labels (user names or dates) and corresponding values (counts). These are then passed to Chart.js to render charts dynamically.

A typical chart configuration used is:

\begin{verbatim}
new Chart(ctx, {
type: 'bar',
data: {
labels: [...],
datasets: [{
label: 'Messages',
data: [...]
}]
}
});
\end{verbatim}

\section{Libraries Used}

\subsection{Python Standard Library}

The following standard libraries were used: \cite{python}

\begin{itemize}
\item \texttt{re} — for regex-based parsing of chat messages and emoji extraction
\item \texttt{datetime} — for handling timestamps and computing time gaps
\item \texttt{json} — for writing structured output to \texttt{data.json}
\item \texttt{sys} — for handling command-line arguments
\end{itemize}

\subsection{External Python Libraries}

\begin{itemize}
\item \texttt{numpy} — used for numerical computations, specifically for computing the activity concentration using vector norms \cite{numpy}
\end{itemize}

\subsection{Frontend Libraries}

\begin{itemize}
\item \textbf{Bootstrap 5} — used for responsive layout and UI components \cite{bootstrap}
\item \textbf{Chart.js} — used for visualizing statistics such as message counts \cite{chartjs}
\end{itemize}



\section{Hurdles Faced}

\begin{itemize}
\item It was difficult for us to design, like coming up with colours, all those styling part took time.
\item Defining ambiguous statistics and making decisions about thresholds
\item Handling time-based calculations correctly
\end{itemize}

\section{Future Improvements}

If given additional time:

\begin{itemize}
\item Add advanced visualizations (graphs)
\item Improve natural language generation
\item Add real-time interaction in frontend
\item Include sentiment analysis
\end{itemize}

\section{Conclusion}

The project successfully demonstrates the integration of data generation, analysis, and visualization. It highlights how behavioral patterns can be extracted from conversational data and presented in an engaging format.

\section{Bibliography}

\begin{thebibliography}{9}

\bibitem{python}
Python Standard Library Documentation.
\url{https://docs.python.org/3/library/}

\bibitem{numpy}
NumPy Documentation.
\url{https://numpy.org/doc/}

\bibitem{bootstrap}
Bootstrap Documentation.
\url{https://getbootstrap.com/docs/5.3/}

\bibitem{chartjs}
Chart.js Documentation.
\url{https://www.chartjs.org/docs/}


\end{thebibliography}


\end{document}
